\section{Discussion}

Currently, although the current GPU implementation does not yet meet the real-time decoding requirement of superconducting platforms, our accelerator design (e.g., tensor update unit and uncertain-bit detector) is well-suited for FPGA or ASIC customization. These units exhibit fixed dataflow patterns and deterministic parallelism, enabling pipelined designs that run at 500 MHz FPGA with reduced latency. After FPGA prototyping, these modules can be further fabricated as ASIC accelerators operating at speeds near 2 GHz, providing a two orders of magnitude improvement in decoding latency. On the other hand, we already satisfy the latency constraints of neutral-atom and ion-trap systems ($\sim 10^2 \mu s$) \cite{wolanski2024ambiguity}. Given that qLDPC codes require long-range connectivity, future qLDPC hardware is likely to be built on neutral-atom platforms \cite{pecorari2025high}. In this case, our decoder can offer both high precision and real-time performance. 

Regarding decoding accuracy, we compared \papername with exact MLD only on very small code (i.e., BB code [[18,4,4]]) because the computational cost of MLD grows exponentially with code size. Even high-performance GPUs (e.g., NVIDIA H100 with 80 GB memory) cannot evaluate medium-sized qLDPC codes with exact MLD due to the limited memory. On BB code [[18,4,4]], \papername achieves an average logical error rate ($3.3\times10^{-4}$) comparable to that of MLD ($3.1\times10^{-4}$), which enables higher noise levels in the future design of large-scale quantum chips. Moreover, higher decoding accuracy also helps relax the qubit connectivity requirements, making the chip less sensitive to underlying hardware constraints \cite{cao2025exact,chubb2021statistical}.

% Looking forward, several stages of the pipeline can benefit from FPGA-based design. For example, the tensor update path in the TCE can be mapped to a fully pipelined contraction unit with a scratchpad hierarchy, removing the general-purpose overheads of tensor cores. Likewise, the uncertain-bit detector of UDE matches FPGA spatial parallelism: each PE can be implemented as a hardwired datapath with fixed matrix selection and reduction logic, enabling deterministic per-bit processing. In addition, a lightweight FPGA control engine can orchestrate refinement rounds without GPU scheduling overhead. With these customized modules, the end-to-end latency can be reduced substantially. For instance, a GPU-bound latency on the order of $10^3 \mu s$ can be decreased to sub-$10^1 \mu s$, yielding both higher throughput and better energy efficiency than a pure GPU design.